% Notice of a Correction.
%
% Article 23 provides that where any two copies differ the master copy
% governs, that a copy found to differ is corrected against it and the
% correction entered, and that the Clerk gives the notice to any citizen who
% asks for it without charge of chits. This is such a notice, for the
% correction carried into this copy at Article 18. See \corr in pact.sty.

\vspace*{0.5in}
\begin{center}
\begin{minipage}{3.6in}
\setlength{\fboxrule}{0.5pt}\setlength{\fboxsep}{14pt}
\fbox{\begin{minipage}{\dimexpr3.6in-29pt\relax}
\ttfamily\footnotesize\raggedright\sloppy
\setlength{\parindent}{0pt}
\setlength{\parskip}{8pt}

\centerline{\scshape\normalsize Notice of a Correction}
\centerline{\scriptsize given under \secref{art:amendment}{sec:amendment:printing-of-a-ratified-amendment}}

\vspace{2pt}

This copy was found to differ from the master copy at
\secref{art:cleaning}{sec:cleaning:days-between-the-request-and-the-departure},
where it read \emph{within two cycles of the request or the sentence} and the
master copy reads \emph{within one cycle}. The copy is corrected against the
master copy at that place. The correction is entered in the archive with the
day of its making.

The Clerk answers for its having been wrong. No citizen is answerable for
having obeyed this copy as it stood, nor is any officer answerable who kept a
citizen the shorter time.

Given to any citizen who asks for it, without charge of chits.

\centerline{\rule{1.4in}{0.4pt}}
\centerline{\scriptsize\rmfamily\scshape Clerk of Judicial}
\end{minipage}}
\end{minipage}

\stampline[2]{Entered \textperiodcentered\ Judicial Archive}

\end{center}
\cleardoublepage
